EE220 ASSIGNMENT ONE - SOLUTIONS
QUESTION ONE
a) Covalent Bond Diagram of Silicon Doped with Pentavalent Impurity
    Si     Si
     \     /
      Si—P—Si
     /     \
    Si     Si
Each Silicon (Si) atom has 4 valence electrons and forms 4 covalent bonds. 
Phosphorus (P) has 5 valence electrons. Four (4) electrons form covalent bonds with four neighbouring silicon atoms.
The fifth electron is not used in bonding and is loosely bound. 
It is therefore released into the conduction band and becomes a free electron.
Pure Silicon doped with a pentavalent impurity becomes an n-type semiconductor.
b) Semiconductor Types and Majority Carriers
Semiconductor Majority Carrier Minority Carrier N-type Electrons Holes P-type Holes Electrons
N-type semiconductor:
· Majority carriers are electrons ·
Produced by doping pure silicon with a pentavalent donor impurity (e.g., phosphorus) · Minority carriers are holes
P-type semiconductor:
· Majority carriers are holes · 
Produced by doping pure silicon with a trivalent acceptor impurity (e.g., boron) · Minority carriers are electrons
c) Doping Process
For John to alter the composition of some intrinsic semiconductors,
he has to deliberately add a small amount of impurity to a pure/intrinsic semiconductor. This process is called doping.
Doping changes the concentration of charge carriers and therefore changes the electrical conductivity of the semiconductor.
· Pentavalent doping: A pentavalent impurity is added leading to extra electrons. This produces an n-type semiconductor. ·
Trivalent doping: A trivalent impurity is added leading to holes. This produces a p-type semiconductor.
Therefore, doping produces majority carriers by adding controlled impurities to an intrinsic semiconductor.
d) Diode Characteristics
i) V-I Characteristic Curve
[Graph should be plotted on graph paper with:
· X-axis: Voltage V (V) from 0 to 0.8V · Y-axis: Current ID (mA) from 0 to 30mA · Plot points from Table 1]
V (V) 0 0.1 0.2 0.3 0.4 0.5 0.6 0.7 0.75 0.8 ID (mA) 0 0 0 0 0 0 0 2.5 5 30
ii) Dynamic (AC) Resistance at Forward Current of 15mA
Using the equation: r_d = \frac{\Delta V_d}{\Delta I_d}
From the graph, at I_d = 15\text{ mA}:
Select two points around 15mA:
· At I_{d1} = 5\text{ mA}, V_{d1} = 0.75\text{ V} · At I_{d2} = 30\text{ mA}, V_{d2} = 0.80\text{ V}
\Delta I_d = 30 - 5 = 25\text{ mA} = 25 \times 10^{-3}\text{ A} \Delta V_d = 0.80 - 0.75 = 0.05\text{ V}
r_d = \frac{0.05}{25 \times 10^{-3}} = 2\ \Omega
iii) Dynamic Resistance Using r_d = \frac{26\text{ mV}}{I_d}
r_d = \frac{26 \times 10^{-3}}{15 \times 10^{-3}} = \frac{26}{15} = 1.73\ \Omega
iv) Comparison of Solutions
Method Value From graph (\Delta V_d/\Delta I_d) 2\ \Omega Using r_d = 26\text{ mV}/I_d 1.73\ \Omega
The two values are close but not identical. The difference arises because:
· The 26\text{ mV}/I_d formula assumes ideal diode behaviour · 
The graphical method uses actual measured data which includes non-ideal effects
v) AC Resistance at 2mA and 20mA
At I_d = 2\text{ mA}: r_d = \frac{26 \times 10^{-3}}{2 \times 10^{-3}} = 13\ \Omega
At I_d = 20\text{ mA}: r_d = \frac{26 \times 10^{-3}}{20 \times 10^{-3}} = 1.3\ \Omega
Comparison and Conclusion:
Current AC Resistance 2 mA 13\ \Omega 20 mA 1.3\ \Omega
Conclusion: As the diode current increases, the AC resistance decreases. 
This is because the diode's dynamic resistance is inversely proportional to the forward current. 
At higher currents, the diode becomes more conductive.
